\section{Related Work}
\label{sec:related}

\paragraph{Retrieved-context memory models.}
The Temporal Context Model \cite{howard2002tcm} and its descendants
\cite{polyn2009cmr,cohen2022cmr3} represent each studied item as a feature
vector bound to a drifting context vector; retrieval reinstates context,
which probabilistically reactivates associated items. CMR3
\cite{cohen2022cmr3} extends CMR with a bi-cell representation of
emotional valence and predicts mood-congruent recall, intrusive
memories, and the effects of emotion dysregulation. These models do not
include a self representation: the context vector drifts with topic, not
identity, and emotional features are content-attached, not agent-attached.

\paragraph{Activation-based memory.}
ACT-R \cite{anderson2007actr} stores declarative chunks with
base-level activation $B_i = \ln \sum_j t_j^{-d}$ summing over past
retrievals. Recent LLM-agent memory architectures port ACT-R-style
activation to dialog memory \cite{memos2025},
combined with surprise or emotional arousal for retention; a recent wave
of 2025 systems carry this style further with cognitive-science-inspired
organisation. None of these treat the self as a distinct generative
entity that gates encoding.

\paragraph{FEP and self.}
Friston's free-energy principle \cite{friston2010fep,friston2017ai}
formalises perception and action as minimising variational free energy
under a generative model. \textit{Self-evidencing} interpretations
\cite{hohwy2016selfevidencing,seth2018beast} cast the agent as actively
sampling sensations that confirm its model of itself. Apps and Tsakiris
\cite{apps2014freeself} apply this to bodily self-recognition; Limanowski
and Blankenburg \cite{limanowski2013minimalself} map minimal phenomenal
selfhood to a hierarchical generative model. The thread we pick up is
the implicit one: if the self is a generative model with precision
$\piself$, that precision should modulate memory the same way it
modulates other generative-model-driven processes.

\paragraph{Self-Reference Effect.}
Rogers, Kuiper, and Kirker \cite{rogers1977sre} showed that
self-referenced encoding produces a robust advantage over semantic,
phonemic, and structural encoding. Symons and Johnson \cite{symons1997meta}
synthesised the meta-analysis. The effect has been a sustained empirical target;
\cite{conway2005smemsys} situates it within the self-memory system,
and \cite{klein2012selfref} reviews conceptual cautions about its
variants and boundary conditions. The empirical pattern is
exceptionally well-documented; what is missing is a computational
mechanism that connects it to a broader theory of mind.

\paragraph{Encoding and forgetting.}
Levels of processing \cite{craik1972lop} and encoding specificity
\cite{tulving1973encodingspec} predict that depth of processing and
context match drive recall. Ebbinghaus \cite{ebbinghaus1885forgetting}
quantified forgetting as a near-power-law of time. Bower
\cite{bower1981mood} established the associative-network account of
mood-congruent memory. Nader and Hardt \cite{nader2011reconsolidation}
reviewed reconsolidation, the phenomenon that retrieved memories become
labile and can be re-encoded. Our model builds on these foundations but
unifies them with the FEP self, which none of the prior work
explicitly formalises in a memory context.

\paragraph{Position of this work.}
We are not claiming the first computational model of the SRE, nor the
first FEP account of self, nor the first emotion-modulated memory model.
We are claiming the first model that wires \emph{all three} together
into a single mechanism whose multiplicative interaction structure
generates falsifiable predictions, and that can be tested as a thin
extension to an existing agent substrate. The closest comparison
target, CMR3, is emotion-only and lacks a self dimension; we leave a
quantitative side-by-side comparison to future work
(\cref{sec:limitations}).
